
\begin{table}[h]
\centering
\caption{Summary of Literature Review}
\label{table:lit}
%\def\arraystretch{1.25}
%\resizebox{\textwidth}{!}{%
\begin{tabular}{|p{3.2cm}|p{3.2cm}|p{3.2cm}|p{3.2cm}|}
\hline
\textbf{Paper Title} & \textbf{Objectives} & \textbf{Methodology} & \textbf{Limitations}  \\ \hline
Machine learning method for cosmetic product recognition: a visual searching approach\cite{1_lit}
 & A recognition system to recognize Cosmetic products with their types, brands, and retailers such that to analyze a customer experience what kind of products and brands they need.                                         & Trained VGG16 and ResNet50 model, SVM classifier & Any extra algorithms for background removal, noise filtering and region of interest for the required object, have not been performed. The images are captured at unconstrained environments with distinct background details, illumination variations, scale invariants within the intra-class images.  \\ \hline
ViSeR: A Visual Search Engine for e-Retail\cite{12_lit}
 & A visual search engine architecture
using deep learning and machine learning techniques, with
a proof-of-concept implementation focusing on the fashion e-retail industry.
                             & CNN model, cosine distance or Manhattan distance, Transfer learning, K-means clustering & Dataset not large enough to employ as the first step in the process of fashion item recommendations.  \\ \hline
Visual E-commerce Application using Deep Learning\cite{13_lit} & A system that scans real-time objects of their choice using an android application and does detection and recommendation. & YOLO model & Scanning image-It is not always possible that the product is physically present with us.   \\ \hline

\end{tabular}
}
\end{table}


